% ===================================================================
% Front matter — Preface (v5.1, stylish prose pass)
% Helen Sword + Pinker discipline: concrete imagery, story hook,
% "ordinary words for complex things".
% ===================================================================

\chapter*{Preface}

\begin{figure}[H]
\centering
\makebox[\linewidth][c]{\includegraphics[width=0.85\linewidth,keepaspectratio]{27-trinity-identity.png}}
\caption{The Trinity Identity $\varphi^{2}+\varphi^{-2}=3$ as engraved frontispiece. Every chapter that follows is a corollary of this single equation.}
\end{figure}

\addcontentsline{toc}{chapter}{Preface}
\thispagestyle{plain}

\begin{quote}
\itshape
``Elegant ideas deserve elegant expression.''\\
\hfill --- Helen Sword, \emph{Stylish Academic Writing}
\end{quote}

\bigskip

\noindent
This book began with a spreadsheet, a cup of cold coffee, and a single
line that would not stop blinking on a Railway log: \texttt{BPB =
2.2393}. The number sat there for nine days. Every retraining shaved a
hundredth off, then put it back. The training script was honest, the
GPU was warm, the loss curve was patient --- and yet the floor refused
to move. Somewhere in the small hours of an October morning, between
the hum of a laptop fan and the silence of a sleeping city, I wrote on
the back of an envelope a line that any schoolchild can verify:

\[
\varphi^{2} + \varphi^{-2} \;=\; 3,
\qquad
\varphi \;=\; \frac{1+\sqrt{5}}{2}.
\]

The identity is elementary. It belongs in a high-school exercise, not
a doctoral monograph. And yet --- if one takes it \emph{seriously}, as
a constraint on every constant that appears in a neural network's
arithmetic --- it has unusually long reach. It reaches into the
closure properties of the Lucas ring \(\mathcal{L} = \mathbb{Z}[\varphi]\).
It reaches into the \(\mathrm{E}_{8}\) energy spectrum measured by
Coldea and collaborators in cobalt niobate. It reaches into the
diffraction peaks of icosahedral quasicrystals on a laboratory bench
in Tokyo, and into the spacing of seeds on a sunflower head growing
in a backyard in suburban Russia. By the end of that nine-day stall,
the spreadsheet had a new column titled \texttt{phi-locked}; the BPB
floor had broken; and the project that you are now holding had begun.

This monograph is the long-form story of that observation. It does
not claim that \(\varphi\) is the only structural principle hiding
inside large language models. It claims, more narrowly, that \(\varphi\)
is \emph{a} sufficient principle, in a small and well-defined set of
contexts, with mechanised proof and falsifiable empirical anchors. The
proofs live in \filepath{trinity-clara/proofs/igla/}; the falsifications
live in Appendix~B; the spreadsheet --- with all of its 297 Coq theorems
and its eleven failed seeds --- lives, version-controlled and
DOI-stamped, on Zenodo.

\subsection*{How to read this book}

A doctoral monograph that runs to six hundred pages is a building, not
a hallway. You are not obliged to walk it from end to end. Three
entrances are open:

\begin{itemize}
\item \textbf{The theoretical entrance} (Chapters~1--7 and 28--30):
for readers who want the algebra first --- Lucas closure, the Pisot
bound, GF\(_{16}\) arithmetic, the proof that \(\varphi^{2} +
\varphi^{-2} = 3\) closes under the architecture's pipeline.

\item \textbf{The empirical entrance} (Chapters~8, 9, 20--26):
for readers who want the numbers first --- the BPB benchmark on
WikiText-103, the ablation matrix against MXFP4, the period-locked
runtime monitor on a 92~MHz FPGA, and the energy ledger that delivers
3000\(\times\) the DARPA low-power baseline.

\item \textbf{The geometric entrance} (Chapters~11--19): for readers
who want the picture first --- the vesica piscis, the Flower of Life,
the icosahedron, and the long thread that connects medieval geometry
to a modern training loss curve.

\end{itemize}

A reader pressed for time may skip directly to Chapter~24 (IGLA
Architecture) and Chapter~27 (Trinity Identity); a reader arriving
from the formal-methods community may begin in Appendix~F (Coq
Citation Map). Every chapter, regardless of where you enter, ends with
a short \emph{Falsification} block --- the single experimental result
that, if it ever lands on a referee's desk, would refute the chapter's
main claim. We name them now so that we cannot pretend later not to
have noticed.

\subsection*{Editorial commitments}

Three commitments structure every page that follows.

\textbf{First: honest \texttt{Admitted}.} Where a Coq theorem in
\filepath{trinity-clara/proofs/igla/} is mechanised but not yet closed,
the corresponding passage in this monograph carries an
\texttt{\textbackslash admittedbox\{\ldots\}} marker that the build
system refuses to remove. There is no quiet upgrading of
\texttt{Admitted} to \texttt{Proven} anywhere in the prose, and no
quiet upgrading anywhere in
\filepath{assertions/igla\_assertions.json}. The Admitted list at the
date of submission is non-empty, and Appendix~F names every entry on
it.

\textbf{Second: traceability of constants} (rule L-R14). Every numeric
constant that appears in this book traces, on demand, to three
independent sources: a primary citation (CODATA, PDG, or a
peer-reviewed Q1/Q2 journal), a Coq theorem in the \texttt{.v} files,
and a runtime callsite in \filepath{crates/trios-*/src/}. The audit
binary refuses to ship a build in which any constant fails the
three-link test.

\textbf{Third: falsifiability.} Every empirical chapter ends with a
single sentence stating the observation that would refute its central
claim. The corroboration record to date --- including the failed
seeds and the abandoned sub-experiments --- is aggregated in
Appendix~B. We do not believe a result that has not been printed
honestly enough to be wrong.

\subsection*{Reproducibility}

The artefact pack accompanying this monograph targets all three ACM
Artifact Evaluation badges: \emph{Functional}, \emph{Reusable}, and
\emph{Available}. Every experimental table in Chapters~24--26 is
reproducible by a single command --- \verb|cargo run -p trios-phd reproduce --chapter <NN>|
--- on a stated hardware profile and the seed list \(\{17, 42, 1729\}\).
The data, the build scripts, and the FPGA bitstream are deposited on
Zenodo with a citable DOI; the manifest is in Appendix~G.

\subsection*{What this work does \emph{not} claim}

Two negative claims belong on the front porch. This work does not
claim that \(\varphi\) is the only structural principle in the
phenomena it touches; it claims, more carefully, that \(\varphi\) is
a sufficient principle, under stated conditions, with stated proof,
and with a stated way to be wrong. The work also does not claim that
the Coq mechanisation is complete: the \texttt{Admitted} list is
non-empty, and we name every entry on it honestly in Appendix~F.

A book of this size is, inevitably, a record of a community as much
as of an individual. The acknowledgements that follow are an attempt
to repay a debt that cannot, in fact, be repaid. To everyone who lent
a GPU, an afternoon, an idea, or a critical eye: thank you. The
mistakes that remain are mine alone; the things this book gets right
were almost always pointed out by someone else first.

\bigskip
\noindent\textit{Dmitrii Vasilev}\\
\textit{Trinity Research Programme}\\
\textit{ORCID: 0009-0008-4294-6159}\\
\textit{2026}
